\textbf{1.\ \textsc{Escape} (8 points)} --- \textit{Päivo Simson.}

During a nuclear weapon test, a bomb is dropped from an airplane at an altitude
of $H = 9\,\mathrm{km}$ and is set to detonate at $h = 500\,\mathrm{m}$ above the
ground. The air drag acting on the bomb during the fall is negligible.
Immediately after releasing the bomb, the plane starts to escape the explosion.
The crew is protected from radiation burst by a protective screen, but the plane
is vulnerable to the shockwave and needs to be as far as possible from the
detonation point.

\textbf{i)} \textit{(1 point)} The maximum speed of the airplane for a straight
level flight (constant altitude) is $v_0$. What is the maximum diving angle so
that the speed does not exceed the speed of sound $c$? The mass of the airplane
is $m$, the air drag force is $F_d = kv^2$, and the gravitational acceleration is
$g$.

For simplicity, assume from now on that the airplane stays at a constant
altitude, flies at a constant speed $v = 190\,\mathrm{m/s}$, and all air
maneuvers are limited by the maximum allowed lift-to-weight ratio $n = 2.5$. The
gravitational acceleration $g = 9.81\,\mathrm{m/s^2}$.

\textbf{ii)} \textit{(1 point)} After releasing the bomb, how much time does the
airplane have before the radiation burst from the bomb hits it?

\textbf{iii)} \textit{(1 point)} What is the smallest possible curvature radius
$R$ of the plane's trajectory and the corresponding bank angle $\alpha$ (cf
figure) of the plane?

\begin{center}\includegraphics[width=0.5\linewidth]{figures/NBPhO_2022_Q1_fig1.png}\end{center}

\textbf{iv)} \textit{(3 points)} Suggest a trajectory for the fastest escape.
Calculate all the parameters that define the shape of the trajectory and its
position relative to the detonation point.

\textbf{v)} \textit{(2 points)} Based on the suggested trajectory, how far is the
airplane from the detonation point when the shockwave hits it? It is estimated
that the safe distance is 25 km. Is the plane able to escape the explosion?
Assume that the average traveling speed of the shock wave is
$u = 350\,\mathrm{m/s}$. For this part, if needed, you can use justified
approximations to simplify the algebra.
